\documentclass[
  aspectratio=169,
]{beamer}

\usepackage[english]{babel}
\usepackage[utf8]{inputenc}
\usepackage[T1]{fontenc}
\usepackage{csquotes}
\usepackage{expl3,biblatex}
\addbibresource{example.bib}
\usepackage{booktabs}
\usetheme{GMU}

\title[Short Presentation Title]{Full Presentation Title}
\subtitle[Short Presentation Subtitle]{Full Presentation Subtitle}
\author[P. J. Fry]{Phillip J. Fry\texorpdfstring{\\}{, }example@gmu.edu}
\institute[GMU Math]{Department of Astronomy, George Mason University}
\date{\today}
\subject{Presentation Subject}
\keywords{the, presentation, keywords}
\begin{document}

\begin{frame}[plain]
\maketitle
\end{frame}

\begin{frame}{Outline}
\tableofcontents
\end{frame}

\section[Short Section 1 Name]{Full Section 1 Name}
\subsection[Short Subsection 1 Name]{Full Subsection 1 Name}

\begin{frame}{Frame Title}{Frame Subtitle}
plain text, \structure{page structure}, \alert{emphasis}
\begin{itemize}
  \item a single-line bullet list item
  \item a bullet list item that is quite long (in order to force a line break),
    which also contains \alert{emphasized text}
  \begin{itemize}
    \item a second-level list item
    \begin{itemize}
      \item a third-level list item
    \end{itemize}
    \item \alert{an emphasized second-level list item}
  \end{itemize}
\end{itemize}
\begin{enumerate}
  \item a numbered list item
  \begin{enumerate}
    \item a second-level list item containing a math expression
      \[ 
        -\mu \Delta \mathbf{u} + (\nabla \mathbf{u}) \mathbf{u} + \nabla p = \mathbf{f} \enskip \text{in } \Omega, \quad \int_{\Omega} p \, dx = 0, \quad u_D \sim \sin(x) \enskip \text{on }\partial \Omega.
      \]
      and a citation \cite{einstein1905tragheit}
  \end{enumerate}
\end{enumerate}
\end{frame}

\subsection[Short Subsection 2 Name]{Full Subsection 2 Name}

\begin{frame}{Text Blocks}
text above a block\footnote{a footnote with an \url{https://address.edu}}
\begin{block}{Block}
  text in a block
\end{block}
\begin{exampleblock}{Example Block}
  text in a block
\end{exampleblock}
\begin{alertblock}{Emphasized Block}
  text in a block
\end{alertblock}
\end{frame}

\subsection[Short Subsection 3 Name]{Full Subsection 3 Name}

\def\age(#1-#2-#3){%
  \the\numexpr(
    \year - #1
    \ifnum\month<#2
      - 1
    \else
      \ifnum\month=#2
        \ifnum\day<#3
          - 1
        \fi
      \fi
    \fi
  )%
}

\begin{frame}{Tables}
\begin{table}
  \begin{tabular}{llc}
    First Name & Surname & Age \\ \midrule
    Albert & Einstein & \age(1879-03-14) \\
    Marie & Curie & \age(1867-11-07) \\
    Thomas & Edison & \age(1847-02-11) \\
  \end{tabular}
  \caption{The great minds of the 19th century}
\end{table}
\end{frame}

\section[Short Section 2 Name]{Full Section 2 Name}
\SectionPage[This is a subtitle]

\begin{frame}[c]{A Numbered Figure}
  \begin{figure}
    \centering
    \includegraphics[width=0.5\textwidth]{branding/gmu-logo-horizontal-black.pdf}
    \caption{This is a caption.}
  \end{figure}
\end{frame}

\begin{frame}[t, plain]{A slide with no footer}
  In case more space is needed at the bottom
\end{frame}

\begin{frame}[plain]
\vfill
\centerline{A plain slide with no title or footer}
\vfill
\end{frame}

\section{\bibname}
\begin{frame}[sideblock]{Bibliography}

\begin{leftblock}
\begin{thebibliography}{9}
  \bibitem{b0}
  P.~J. Fry, T. Leela and H.~J. Farnsworth,
  \emph{Optimal delivery routes through the Bermuda Tetrahedron: a time-dependent shortest-path problem with anchovy constraints},
  \newblock{J. Interplanet. Delivery Sci. {\bf 42} (3001), 1--1000}.
  
  \bibitem{b1}
  J.~A. Zoidberg and P.~J. Fry,
  \emph{On the absence of the delta brainwave in a twentieth-century human: a case study in accidental immunity to telepathic attack},
  \newblock{Decapod. J. Hum. Med. {\bf 7} (3002), 12--13}.
\end{thebibliography}
\end{leftblock}%
\begin{rightblock}{Acknowledgements}
Or funding, contact\\
or any other information.\\[1em]
example@gmu.edu
\end{rightblock}

\end{frame}

\makeoutro

\end{document}